\documentclass[11pt]{article}
\usepackage[T1]{fontenc}
\usepackage[utf8]{inputenc}
\usepackage[a4paper,margin=25mm]{geometry}
\usepackage{lmodern,microtype,hyperref}
\hypersetup{colorlinks=true,urlcolor=blue}
\title{Prospectus closure: language and option selection}
\author{LegalMath development note}
\date{5 October 2026}
\begin{document}
\maketitle
BASF's March 2032 notes illustrate a problem that downloading a prospectus alone
cannot solve. The final terms select fixed-rate Option I from the September 2022
programme, specify German as the controlling language, and delete alternatives
that are unselected, incomplete or struck out. A reader must apply those choices
before interpreting the template's clauses.\footnote{BASF final terms,
6 March 2023, pp.~2--3 and 7; \href{../../../../prospectus/evidence-master/requests/044-basf-final/response.body}{retained original}.}

The new source record binds the final terms to the base prospectus of
9 September 2022 and its first supplement of 27 February 2023. It checks
21 explicit selections. Page~6 selects an issuer call during
8 December 2031--7 March 2032 but marks a separately worded call tied to interest
payment dates as unavailable. A search for issuer calls cannot distinguish those
provisions. Likewise, the annual ICMA choice excludes short or long coupons,
although other day-count alternatives remain printed in the programme.

The controlling German text preserves further qualifications. The ranking clause
recognises priority created by mandatory law. The payment-day clause denies
additional interest for the specified delay. Section~15 expressly states that
the English translation is nonbinding.\footnote{BASF base prospectus,
pp.~110, 112, 114 and 133; \href{../../../../prospectus/evidence-closure/requests/037-basf-base-september-2022-exchange/response.bin}{retained original}.}
These passages support bounded source decisions. Other optional clauses,
incorporated financial statements and applicable agreements still require review.

\newpage
Deutsche Bank's 2025 AT1 prospectus presents a different obstacle. On pages~36--38,
German occupies the left column and English the right. The earlier extraction
preserved visual rows, causing the reader to join fragments from both languages.
The reviewed derivative separates the columns using word coordinates and retains
every word in its language or page-number region.\footnote{Deutsche Bank,
AT1 Notes Prospectus, ISIN DE000A460DG7, PDF pp.~36--38;
\href{../../../../prospectus/classification-additions/originals/deutsche-at1-2025.pdf}{retained original}.}

The distinction matters in section~2(7). The resolution authority may write
claims down to zero, or convert them into ordinary shares of the issuer, a group
entity or a bridge bank, or into other ownership instruments qualifying as
Common Equity Tier~1. The last alternative prevents a conclusion that conversion
must be exclusively into common shares. The paragraph also preserves the
no-event-of-default qualification. Those words are now available as a contiguous
English passage, alongside the retained German column.

\medskip
123 engineering tests passed. The source-admission tests reject lost negation,
ambiguous selections, inconsistent bilingual fields, changed originals or images,
overlapping geometry and derivatives that no longer reconstruct from the reviewed
words. The earlier OCR and conditional financial checks also pass.

The frozen reader generated 3 records, including 1 unresolved,
on the original three pages and 3 records, including 1 unresolved,
after column separation. These figures describe changed segmentation and
classification on known development examples. They do not measure legal accuracy.
The engineering result is narrower: the reviewed words and material qualifiers
survive the optional extraction path, and provenance checks reject alteration.
A source-specific annotation now records the disclosed write-down power and both
conversion alternatives. It abstains on unreviewed rewording or added exceptions;
actual legal applicability and mandatory common-share-only conversion remain
unknown. This exact reviewed construction is not a general semantic rule.

Independent legal adjudication, further page qualification, complete operative
documents and actual issuer or investor facts remain necessary. No production
default changed, and this round made no new network request. Historical evidence
and existing independent-review forms remain intact. The next source task is to
finish BASF's German clause selection; the next extraction task is to qualify
the remaining Deutsche pages before widening semantic repair.
\end{document}
